\title{Controlling Text-to-Image Diffusion by Orthogonal Finetuning}

\begin{abstract}
Large-scale text-to-image diffusion models exhibit remarkable proficiency in producing photorealistic images based on text prompts. A key challenge that arises is how to effectively steer or adapt these advanced models for a variety of downstream applications. To address this, we propose a theoretically grounded finetuning approach called Orthogonal Finetuning (OFT), designed for tailoring text-to-image diffusion models to new tasks. Distinct from prior finetuning strategies, OFT guarantees the preservation of hyperspherical energy—a metric capturing the pairwise relationships between neurons when projected onto a unit hypersphere. Our analysis shows that maintaining this property is essential for retaining the models' original semantic generation capabilities. To further enhance the stability of the finetuning process, we introduce Constrained Orthogonal Finetuning (COFT), which augments OFT by enforcing an explicit radius constraint on the hypersphere. We focus our investigation on two central finetuning tasks in the text-to-image domain: (1) subject-driven generation, where the model is required to synthesize images of a given subject (from only a handful of provided images) placed within novel contexts using text prompts, and (2) controllable generation, in which the model is provided with supplementary control information. Through extensive experiments, we demonstrate that OFT delivers superior image generation quality and faster convergence compared to existing approaches.
\end{abstract}

\section{Introduction}

Contemporary text-to-image diffusion models~\cite{saharia2022photorealistic,ramesh2022hierarchical,rombach2022high} have set a high bar in terms of fidelity and flexibility when synthesizing images from textual descriptions. However, the reliance on text alone often leads to ambiguities, limiting precise and fine-level control over the generated outputs. In this work, we focus on two primary text-to-image tasks that exemplify these limitations:

\begin{itemize}[leftmargin=*,nosep]

    \item \textbf{Subject-driven generation}~\cite{ruiz2023dreambooth}: Starting from only a limited set of images depicting a subject, the goal is to produce new images featuring the same subject situated in diverse settings, guided by a textual prompt.
    \item \textbf{Controllable generation}~\cite{zhang2023adding,mou2023t2i}: Here, the objective is to incorporate additional control cues (such as canny edge maps or segmentation masks) alongside text, and generate images that respect both the control input and the descriptive prompt.
\end{itemize}

At their core, both challenges revolve around the problem of efficiently adapting text-to-image diffusion models while maintaining the generative strengths acquired during pretraining. An ideal finetuning strategy should fulfill the following criteria: (1) \emph{training efficiency}, entailing reduced numbers of trainable parameters and minimized training epochs, and (2) \emph{generalizability preservation}, or the ability to sustain high-quality and varied generation outputs. Current finetuning approaches predominantly employ either direct modification of neuron weights using a low learning rate (\emph{e.g.}, \cite{ruiz2023dreambooth}) or integration of a compact, parameterized component that alters certain weights (\emph{e.g.}, \cite{hulora2022,zhang2023adding}). Although these methods are straightforward, they do not inherently ensure that generative performance from pretraining is retained. Additionally, a systematic framework for selecting an optimal finetuning method or for tuning key hyperparameters, such as training epochs, remains undeveloped. The core obstacle lies in the absence of an objective metric to evaluate how well a finetuned model retains pretrained generative competencies. Prevailing finetuning techniques operate under the assumption that a smaller Euclidean distance between finetuned and pretrained model parameters equates to better retention of pretrained features, often leading to the adoption of very small learning rates. However, we contend that proximity in Euclidean space does not necessarily equate to faithful preservation of semantics. Consequently, developing a more structurally meaningful metric to distinguish between finetuned and original models could significantly enhance both the maintenance of pretraining performance and the overall reliability of the finetuning process.

Drawing from prior empirical evidence suggesting that hyperspherical similarity effectively encapsulates semantic relationships~\cite{liu2017deep,liu2018decoupled,chen2020angular}, our approach leverages hyperspherical energy~\cite{liu2018learning} to describe the pairwise associations between neurons within a layer. Hyperspherical energy aggregates the hyperspherical similarity (such as cosine similarity) among all neuron pairs at the same layer, providing a quantitative measure of neuron distribution uniformity on the unit hypersphere~\cite{liu2021learning}. We posit that an optimal finetuned model should maintain a hyperspherical energy value closely aligned with that of its pretrained counterpart. Simply introducing a regularization term to keep hyperspherical energy constant during finetuning does not inherently guarantee minimization of energy discrepancy. To address this, we exploit an invariance property: hyperspherical energy remains unchanged if an identical orthogonal transformation is applied to all neurons. Based on this principle, we introduce Orthogonal Finetuning~(OFT), a method for adapting large text-to-image diffusion models to new tasks while preserving their hyperspherical energy. OFT’s core mechanism involves learning a single orthogonal transformation for each layer’s neurons to conserve pairwise angles, effectively shifting their basis in the canonical coordinate system. Recognizing that reduced Euclidean distance between the parameters of the finetuned and pretrained models correlates with greater retention of pretrained capabilities, we further develop Constrained Orthogonal Finetuning~(COFT)—a variant that restricts the finetuned parameters to lie on a hypersphere of fixed radius centered at the pretrained neurons.

The rationale underlying the effectiveness of orthogonal transformation in finetuning neurons draws in part from the 2D Fourier transform, which enables the decomposition of images into magnitude and phase spectra. The phase spectrum—encapsulating the angular relationship between input and basis—retains most of the semantic content. Notably, if the phase spectrum of an image is combined with an arbitrary magnitude spectrum, the reconstructed image preserves its semantic integrity. This observation highlights that adjusting neuron orientations plays a crucial role in altering the semantics of the generated output, a central objective in both subject-driven and controllable image generation. Nonetheless, permitting excessive variation in neuron directions can compromise the generative capabilities established during pretraining. To address this, we advocate for preserving the angles between all neuron pairs, predicated on the notion that these angular relationships are fundamental to neural network knowledge representation. Consequently, our approach involves learning a shared orthogonal transformation within each layer so that the hyperspherical energy remains constant.

Our methodology is also informed by orthogonal over-parameterized training~\cite{liu2021orthogonal}, a technique for training classification neural networks from a randomly initialized state via orthogonal transformations. Since such initialization tends to yield low expected hyperspherical energy, the method in~\cite{liu2021orthogonal} seeks to maintain this low energy throughout training—a property that has been linked to improved generalization in classification tasks~\cite{liu2018learning,lin2020regularizing}. The flexibility provided by orthogonal transformation is shown in~\cite{liu2021orthogonal} to suffice for developing robust classifiers. In contrast, our work centers on finetuning text-to-image diffusion models to enhance controllability and bolster generative quality. There are two principal distinctions between OFT and~\cite{liu2021orthogonal}. First, while~\cite{liu2021orthogonal} is constructed to minimize hyperspherical energy, OFT strives to maintain the energy level of the pretrained model, thereby preserving its essential structure during finetuning—an important consideration since further energy minimization could disrupt existing semantic frameworks in diffusion models. Second, OFT is formulated to minimize the departure from the pretrained parameters, leading to a constrained approach, whereas~\cite{liu2021orthogonal} does not impose this form of regularization. Ultimately, finetuning hinges on balancing adaptability with retention of pre-existing capabilities; we contend that OFT is well-suited to achieve this equilibrium. Our main contributions are summarized as follows:

\begin{itemize}[leftmargin=*,nosep]

    \item We introduce a new finetuning technique, Orthogonal Finetuning (OFT), designed to enhance the controllability of text-to-image diffusion models. To further promote robustness during finetuning, we also devise a constrained variant that restricts the model’s angular divergence from the original pretrained parameters.
    \item In contrast to conventional finetuning strategies, OFT carries out parameter updates while provably maintaining the model’s hyperspherical energy. Our experiments suggest that preserving this quantity is crucial for retaining the pretrained model’s generative semantic fidelity.
    \item We employ OFT in two principal application areas: subject-driven generation and controllable generation. Through a rigorous experimental evaluation, we report notable improvements over existing approaches in generation fidelity, training stability, and convergence speed. Additionally, OFT demonstrates higher sample efficiency, successfully adapting with fewer images and training iterations.
    \item For the evaluation of controllable generation, we present a novel control consistency metric, which quantifies controllability by comparing the estimated control signal inferred from a generated image against the original conditioning input. This metric further substantiates OFT’s superior control capability.
\end{itemize}

\section{Related Work}

\textbf{Text-to-image diffusion models}. The field of text-to-image synthesis has witnessed substantial advancements~\cite{saharia2022photorealistic,ramesh2022hierarchical,rombach2022high,gu2022vector,nichol2022glide}, propelled largely by developments in diffusion-based generative modeling~\cite{ho2020denoising,song2021score,song2021denoising,dhariwal2021diffusion} and vision-language pretraining~\cite{su2020vlbert,lu2019vilbert,radford2021learning,wang2021simvlm,li2022blip,li2023blip,alayrac2022flamingo,schuhmann2022laion}. In particular, GLIDE~\cite{nichol2022glide} and Imagen~\cite{saharia2022photorealistic} implement diffusion in the pixel domain, with GLIDE jointly training its text encoder and diffusion backbone on text-image pairs, whereas Imagen utilizes a frozen pretrained language model as its encoder. Latent space-based training is adopted by both Stable Diffusion~\cite{rombach2022high} and DALL-E2~\cite{ramesh2022hierarchical}: Stable Diffusion leverages VQ-GAN~\cite{esser2021taming} for learning its visual latent codes, while DALL-E2 uses CLIP~\cite{radford2021learning} for constructing a shared multimodal embedding space for both modalities. Beyond diffusion-based methods, research into text-to-image generation has explored generative adversarial networks~\cite{reed2016generative,zhang2017stackgan,xu2018attngan,li2019controllable} and autoregressive architectures~\cite{ramesh2021zero,ding2021cogview,wu2022nuwa,yuscaling}. As a general-purpose finetuning scheme, OFT can be integrated with any text-to-image diffusion model regardless of architectural choices.

\textbf{Subject-driven generation}. To address subject preservation, approaches such as~\cite{avrahami2022blended,nichol2022glide} incorporate user-supplied masks to constrain changes during editing. Alternatively, inversion-based frameworks~\cite{choi2021ilvr,dhariwal2021diffusion,ramesh2022hierarchical,gal2022image} permit background or contextual modifications while attempting to maintain the subject intact. Recent techniques like~\cite{hertz2022prompt} facilitate both localized and holistic image edits without requiring explicit input masks. However, most methods mentioned above have limited ability to preserve fine identity-specific features. The Pivotal Tuning method~\cite{roich2022pivotal} addresses identity retention by regularizing the generator around a latent code obtained via inversion, while~\cite{nitzan2022mystyle} constructs a personalized face prior from a gallery of individual portraits. Although~\cite{casanova2021instance} allows for generating various instances, it does not reliably safeguard unique instance details. DreamBooth~\cite{ruiz2023dreambooth} fine-tunes a diffusion model by incorporating custom tokens and a handful of subject images, optimizing a reconstruction objective along with a class-prior loss to encourage subject fidelity. Our OFT framework is compatible with the DreamBooth pipeline; however, instead of conventional gradient-based tuning with a small learning rate, OFT modifies model parameters via orthogonal transformations.

\textbf{Controllable generation}. Image-to-image translation can be conceptualized as a problem of controllable generative modeling. Earlier approaches largely rely on conditional generative adversarial networks~\cite{isola2017image,zhu2017toward,wang2018high,park2019semantic,choi2018stargan} to address this challenge. More recently, diffusion models have emerged as a powerful tool for image-to-image translation~\cite{saharia2022palette,voynov2022sketch,wang2022pretraining}. Advances such as ControlNet~\cite{zhang2023adding} have enabled precise manipulation of pretrained diffusion models by finetuning them with additional control signals, yielding strong results in controllable generation. Another notable approach, T2I-Adapter~\cite{mou2023t2i}, similarly adapts pretrained diffusion models via finetuning to improve control over generated outputs. Building on the task formulation in \cite{zhang2023adding,mou2023t2i}, we employ OFT for finetuning pretrained diffusion models, which leads to enhanced controllability with reduced training data requirements and fewer finetuned parameters. Importantly, OFT incurs no extra computational cost during test-time inference.

\textbf{Model finetuning}. Adapting large-scale pretrained models for downstream tasks has gained widespread adoption in recent years~\cite{devlin2018bert,brown2020language,he2022masked}. Within this paradigm, adaptation techniques (\emph{e.g.}, \cite{hulora2022,houlsby2019parameter,pfeiffer2020adapterfusion}) are especially prevalent in the natural language processing domain. The most closely related approach to OFT is LoRA~\cite{hulora2022}, which leverages low-rank parameter updates for additive modifications during finetuning. By contrast, OFT applies a multiplicative, layer-wide orthogonal transformation to neuron weights. This operation mathematically guarantees preservation of pairwise angular relationships among neurons within each layer, resulting in improved finetuning stability.

\section{Orthogonal Finetuning}

\subsection{Why Does Orthogonal Transformation Make Sense?}

We begin our analysis by exploring the motivation for employing orthogonal transformations in finetuning text-to-image diffusion models. To clarify this reasoning, we divide the question into two parts: (1) the rationale for finetuning neuron angles (i.e., their directions), and (2) the justification for using orthogonal transformations specifically for angle finetuning.

To address the first question, we build upon the empirical findings in \cite{liu2018decoupled,chen2020angular}, which highlight that differences in angular features effectively represent the semantic disparity between representations. SphereNet~\cite{liu2017deep} demonstrates that confining all neural network activations to a unit hypersphere does not reduce model capacity and often results in improved generalization, suggesting that neuron orientations predominantly encode crucial data attributes. To further illustrate the significance of neuronal directionality, we conduct a simple illustrative experiment as depicted in Figure~\ref{fig:toy_ae}. Here, a conventional convolutional autoencoder is trained on images of flowers. During training, the feature map is computed with the conventional inner product (where $z$ is the convolution between the kernel $\bm{w}$ and the input patch $\bm{x}$). At inference, we compare three methods of deriving the feature map: (a) replicating the inner product from training, (b) using only the magnitude, and (c) employing only the angle (directional component). Figure~\ref{fig:toy_ae} reveals that representations relying solely on angular components can nearly reconstruct the original images, whereas those using only the magnitude retain no informative content. It is important to note that cosine-based activations (c) are not used during training; all learning proceeds via the inner product. These observations strongly suggest that neuron orientation predominantly encodes semantic content within images, indicating that adjusting neuron direction should be prioritized for manipulating semantic information.

Regarding the second question, a natural method to adjust neuron directions is to simultaneously apply a rotation or reflection to all neurons within a given layer, inherently invoking orthogonal transformations. While it might be possible to increase flexibility by introducing angle-dependent transformations for different neurons, we find that orthogonal transformations strike a practical balance between adaptability and structure. Additionally, as argued by \cite{liu2021orthogonal}, orthogonal transformations are adequately expressive for neural network optimization. To empirically support this claim, we perform a controlled generation experiment following the ControlNet~\cite{zhang2023adding} configuration, with results presented in Figure~\ref{fig:ortho}. The standard OFT (Orthogonal Finetuning Transformation) shows consistent performance, achieving precise manipulation capabilities by epoch 20. In contrast, removing the orthogonality constraint from OFT results in unrealistic outputs and fails to achieve any meaningful control. This experiment underscores the crucial role played by the orthogonality constraint within OFT.

\subsection{General Framework}

Traditional finetuning methods often employ gradient descent with a low learning rate to update either the entire model or select layers. The modest learning rate serves to implicitly constrain the deviation from the original pretrained weights, so conventional finetuning effectively performs optimization with a hidden regularization term on $\thickmuskip=2mu \medmuskip=2mu \| \bm{M} - \bm{M}^0\|$, where $\bm{M}$ are the finetuned parameters and $\bm{M}^0$ are the original pretrained weights. While this type of implicit regularization seems reasonable, it can still allow too much freedom when adapting large models. To counter this, LoRA proposes a supplementary restriction by limiting the rank of the parameter update, specifically enforcing $\thickmuskip=2mu \medmuskip=2mu \text{rank}(\bm{M} - \bm{M}^0)=r'$ for some small $r'$. By contrast, OFT imposes a constraint focused on preserving the pairwise neuron relationships through hyperspherical energy: $\thickmuskip=2mu \medmuskip=2mu \| \text{HE}(\bm{M})-\text{HE}(\bm{M}^0) \|=0$, where $\text{HE}(\cdot)$ indicates the hyperspherical energy of a weight matrix. 

To illustrate, consider a dense layer $\thickmuskip=2mu \medmuskip=2mu \bm{W}=\{\bm{w}_1,\cdots,\bm{w}_n\}\in\mathbb{R}^{d\times n}$ where each $\bm{w}_i\in\mathbb{R}^{d}$ is an individual neuron (with $\bm{W}^0$ as the pretrained weights). The layer computes the output vector $\thickmuskip=2mu \medmuskip=2mu \bm{z}\in\mathbb{R}^n$ as $\thickmuskip=2mu \medmuskip=2mu \bm{z}=\bm{W}^\top\bm{x}$ given the input $\thickmuskip=2mu \medmuskip=2mu \bm{x}\in\mathbb{R}^d$. The principle behind OFT is to minimize the gap in hyperspherical energy between the adapted and original weights:

\begin{equation}\label{eq:oft_principle}
\footnotesize
\min_{\bm{W}} \left\| \text{HE}(\bm{W})-\text{HE}(\bm{W}^0)\right\|~~\Leftrightarrow~~\min_{\bm{W}} \bigg{\|} \sum_{i\neq j} \|\hat{\bm{w}}_i-\hat{\bm{w}}_j\|^{-1}-\sum_{i\neq j} \|\hat{\bm{w}}^0_i-\hat{\bm{w}}^0_j\|^{-1}\bigg{\|}
\end{equation}

Here, $\thickmuskip=2mu \medmuskip=2mu \hat{\bm{w}}_i:=\bm{w}_i/\|\bm{w}_i\|$ represents the normalized $i$-th neuron, and the hyperspherical energy for $\bm{W}$ is given by $\thickmuskip=2mu \medmuskip=2mu \text{HE}(\bm{W}):= \sum_{i\neq j} \|\hat{\bm{w}}_i-\hat{\bm{w}}_j\|^{-1}$. Notably, equation~\eqref{eq:oft_principle} attains a minimum value of zero, which is achievable if and only if $\bm{W}$ is a rotated or reflected version of $\bm{W}^0$, i.e., $\thickmuskip=2mu \medmuskip=2mu \bm{W}=\bm{R}\bm{W}^0$ for some orthogonal matrix $\thickmuskip=2mu \medmuskip=2mu \bm{R}\in\mathbb{R}^{d\times d}$ (where the matrix determinant being $1$ or $-1$ corresponds to a rotation or reflection, respectively). OFT leverages this insight, restricting finetuning to learning a shared orthogonal transformation at each layer. More precisely, for a pretrained dense layer $\thickmuskip=2mu \medmuskip=2mu

\begin{equation}\label{eq:oft_general}
    \footnotesize
    \bm{z} = \bm{W}^\top\bm{x} = (\bm{R} \cdot \bm{W}^0)^\top\bm{x},~~~\text{s.t.}~\bm{R}^\top\bm{R} = \bm{R}\bm{R}^\top = \bm{I}
\end{equation}

Here, $\bm{W}$ stands for the weight matrix refined by OFT, while $\bm{I}$ represents the identity matrix. An overview of OFT is presented in Figure~\ref{fig:block}. Drawing a parallel to the zero-initialization strategy utilized by LoRA, it is necessary that OFT enables fine-tuning to proceed from the original pretrained parameters $\bm{W}^0$. To ensure this, we initialize the orthogonal matrix $\bm{R}$ to the identity matrix, thereby ensuring that the fine-tuned parameters initially coincide with the pretrained ones. For maintaining the orthogonality constraint on $\bm{R}$ throughout optimization, we employ differentiable orthogonalization techniques referenced in \cite{liu2021orthogonal,lezcano2019cheap}. Methods for achieving computationally efficient orthogonality will be explored in the following section.

\subsection{Efficient Orthogonal Parameterization}\label{sect:eff_ortho}

Classical orthogonalization algorithms, such as the Gram-Schmidt process, though differentiable, tend to impose considerable computational cost in practice~\cite{liu2021orthogonal}. To enhance computational efficiency, we utilize the Cayley transform to obtain the orthogonal matrix. In particular, the orthogonal matrix is constructed by setting $\thickmuskip=2mu \medmuskip=2mu \bm{R} = (\bm{I} + \bm{Q})(I - \bm{Q})^{-1}$, where $\bm{Q}$ is a skew-symmetric matrix so that $\thickmuskip=2mu \medmuskip=2mu \bm{Q} = -\bm{Q}^\top$. This approach trades a minor limitation for computational gains—the Cayley transform can only yield orthogonal matrices with determinant equal to $1$, i.e., belonging to the special orthogonal group. Empirically, we observe that this restriction has negligible impact on model performance. However, even with the Cayley parameterization, representing a large orthogonal matrix $\bm{R}$ for high-dimensional cases (large $d$) can remain parameter inefficient. To reduce parameter count, we introduce a block-diagonal structure for $\bm{R}$, partitioning it into $r$ independent blocks, leading to the following representation:

\begin{equation}
    \footnotesize
    \bm{R} = \text{diag}(\bm{R}_1, \bm{R}_2, \cdots, \bm{R}_r) = \begin{bmatrix}
    \bm{R}_1 \in \text{O}(\frac{d}{r}) &  &\\
    &\ddots & \\
    & & \bm{R}_r \in \text{O}(\frac{d}{r})
    \end{bmatrix} \in \text{O}(d)
\end{equation}

Here, $\text{O}(d)$ represents the orthogonal group in $d$ dimensions, and both $\thickmuskip=2mu \medmuskip=2mu \bm{R}\in\mathbb{R}^{d\times d}$ and $\thickmuskip=2mu \medmuskip=2mu \bm{R}_i\in\mathbb{R}^{d/r\times d/r},\forall i$ are orthogonal matrices. When $\thickmuskip=2mu \medmuskip=2mu r=1$, the block-diagonal orthogonal matrix reduces to the standard full orthogonal matrix without additional constraints. For a square orthogonal matrix of dimension $d\times d$, the total parameter count is $\thickmuskip=2mu \medmuskip=2mu d(d-1)/2$, giving a computational complexity of $\mathcal{O}(d^2)$. In contrast, an orthogonal matrix with an $r$-block diagonal structure possesses $\thickmuskip=2mu \medmuskip=2mu d(d/r-1)/2$ free parameters, yielding a complexity of $\thickmuskip=2mu \medmuskip=2mu \mathcal{O}(d^2/r)$. Additionally, parameter sharing across all blocks can be adopted, meaning $\thickmuskip=2mu \medmuskip=2mu\bm{R}_i=\bm{R}_j$ for every $i\neq j$, to reduce the parameter complexity further to $\mathcal{O}(d^2/r^2)$. It is worth emphasizing that, irrespective of such parameter reductions, $\bm{R}$ remains an orthogonal matrix, thus retaining the full hyperspherical energy.

We further compare the parameter efficiency of OFT with that of LoRA. For LoRA with a low-rank parameter $r'$, the trainable parameter count is $\thickmuskip=2mu \medmuskip=2mu r'(d+n)$. If both $r$ and $r'$ are taken as fractions of the model dimension $d$, for example, by setting $\thickmuskip=2mu \medmuskip=2mu r=r'=\alpha d$ with $\thickmuskip=2mu \medmuskip=2mu 0<\alpha\leq 1$, then LoRA has complexity $\mathcal{O}(d^2+dn)$, while OFT requires only $\mathcal{O}(d)$. To clarify the difference, consider a model where the weight matrix has dimensions $\thickmuskip=2mu \medmuskip=2mu 128\times 128$; in this scenario, LoRA results in $2,048$ trainable parameters for $\thickmuskip=2mu \medmuskip=2mu r'=8$, whereas OFT requires $960$ trainable parameters for $\thickmuskip=2mu \medmuskip=2mu r=8$ (assuming block sharing is not applied).

\subsection{Constrained Orthogonal Finetuning}

To further restrict the adaptability of the conventional OFT approach, one may enforce that the finetuned parameters remain in close proximity to the pretrained model. Specifically, in COFT, the forward computation is given by:

\begin{equation}\label{eq:coft_general}
    \footnotesize
    \bm{z}=\bm{W}^\top\bm{x}=(\bm{R}\cdot\bm{W}^0)^\top\bm{x},~~~\text{s.t.}~\bm{R}^\top\bm{R}=\bm{R}\bm{R}^\top=\bm{I},~\norm{\bm{R}-\bm{I}}\leq \epsilon
\end{equation}

This imposes two constraints: an orthogonality condition and a bound on the $\epsilon$-deviation from the identity matrix. While the orthogonality can be ensured using the Cayley parameterization as discussed in Section~\ref{sect:eff_ortho}, integrating the $\epsilon$-deviation constraint directly into a Cayley-parameterized matrix is challenging. To provide intuition into the Cayley transform, we use the Neumann series to approximate $\thickmuskip=2mu \medmuskip=2mu \bm{R}=(\bm{I}+\bm{Q})(I-\bm{Q})^{-1}$ by $\thickmuskip=2mu \medmuskip=2mu \bm{R

series converges under the operator norm. As a result, we are able to incorporate the constraint $\thickmuskip=2mu \medmuskip=2mu\|\bm{R}-\bm{I}\|\leq\epsilon$ within the Cayley transform itself. This adjustment translates the original condition to $\thickmuskip=2mu \medmuskip=2mu \|\bm{Q}-\bm{0}\|\leq\epsilon'$, with $\bm{0}$ representing the zero matrix and $\epsilon'$ signifying a different, specifically chosen tolerance (distinct from $\epsilon$). This revised condition on $\bm{Q}$ is straightforward to impose using projected gradient descent. For identity initialization of the orthogonal matrix $\bm{R}$, $\bm{Q}$ is set initially to the zero matrix. Conceptually, COFT combines two explicit regularizations: minimal change in hyperspherical energy and a strict limitation on how far the model deviates from its original, pretrained configuration. While conventional finetuning schemes usually incorporate the second constraint in an implicit manner, COFT formalizes it. We have observed that, although OFT exhibits strong empirical results, the introduction of COFT's explicit constraint on deviation further enhances the stability of the finetuning process. In Figure~\ref{fig:eps_coft}, we illustrate the effect of varying $\epsilon$ on COFT performance. This parameter regulates the degree of permissible adaptation during finetuning: increasing $\epsilon$ yields behavior increasingly akin to OFT, whereas decreasing $\epsilon$ draws the COFT solution closer to that of the initial pretrained diffusion model.

\subsection{Re-scaled Orthogonal Finetuning}

We introduce a straightforward modification to OFT by additionally endowing each neuron with an independently learned scaling factor. Our motivation stems from the observation that scaling does not alter the hyperspherical energy due to subsequent normalization. The revised forward operation becomes: $\thickmuskip=2mu \medmuskip=2mu\bm{z}=(\bm{R}\bm{W}^0\bm{D})^\top\bm{x}$\footnote{\textbf{Errata}: The NeurIPS camera-ready version incorrectly describes the forward computation for re-scaled OFT as $\thickmuskip=2mu \medmuskip=2mu\bm{z}=(\bm{D}\bm{R}\bm{W}^0)^\top\bm{x}$. However, our actual implementation was correct, so the results reported in Appendix~\ref{app:rescaled} remain valid.}, where $\thickmuskip=2mu \medmuskip=2mu \bm{D}=\text{diag}(s_1,\cdots,s_n)\in\mathbb{R}^{n\times n}$ is a diagonal matrix of trainable, strictly positive scalars $s_1,\ldots,s_n$. Contrary to the original OFT formulation in Eq.~\eqref{eq:oft_general}, where learning is limited to $\bm{R}$, the re-scaled OFT jointly optimizes both $\bm{D}$ and $\bm{R}$. This adjustment offers greater adaptation capacity to OFT, yet introduces only a minimal parameter overhead. For the core experiments, we adhere to vanilla OFT to isolate the benefits of orthogonal transformations alone; nevertheless, our findings indicate that the re-scaled OFT variant typically provides additional advantages (see Appendix~\ref{app:rescaled} for empirical support).

\section{Intriguing Insights and Discussions}

\textbf{OFT accommodates diverse architectures}. OFT, by its design, can be seamlessly incorporated into any neural network architecture. Within Transformers, it is common practice to apply LoRA to the attention layers~\cite{hulora2022}. For an equitable comparison with LoRA, our experiments limit the application of OFT to the attention weight matrices. Beyond standard fully connected layers, OFT is also highly compatible with convolutional layers. The block-diagonal structure of $\bm{R}$ leads to notable effects in convolutional contexts (in contrast with LoRA). Specifically, if the block count equals the number of input channels, each block operates on a distinct channel—effectively mimicking depth-wise convolution kernels~\cite{chollet2017xception}. Alternatively, sharing blocks across channels allows the same orthogonal transformation to be applied to all neuron channels. We offer a preliminary investigation into applying OFT to convolutional layers in Appendix~\ref{app:convolution}.

\textbf{Connection to LoRA}. LoRA addresses the issue of preserving pretrained knowledge by adding a low-rank matrix to the existing weight matrix, thus limiting drastic shifts in its information content. In comparison, OFT preserves pretrained information by constraining the transformation applied to the pretrained weight matrix to be orthogonal (and full-rank), ensuring that pretraining representations are not compromised. The forward pass in OFT can be recast as $\thickmuskip=2mu \medmuskip=2mu \bm{z}=(\bm{R}\bm{W}^0)^\top\bm{x}=(\bm{W}^0+(\bm{R}-\bm{I})\bm{W}^0)^\top\bm{x}$, where $\thickmuskip=2mu \medmuskip=2mu (\bm{R}-\bm{I})\bm{W}^0$ serves a similar role to the low-rank update in LoRA. Because $\bm{W}^0$ is generally full-rank, OFT performs a low-rank update to the weight matrix if $\thickmuskip=2mu \medmuskip=2mu \bm{R}-\bm{I}$ is low-rank. While LoRA utilizes a rank parameter $r'$ to control trainable parameters, OFT introduces a diagonal block parameter $r$ for the same purpose. Fundamentally, LoRA and OFT present complementary strategies for achieving parameter efficiency: LoRA relies on exploiting low-rank structures, whereas OFT leverages sparsity via block-diagonal orthogonal constraints.

\textbf{Why OFT converges faster}. On one side, Figure~\ref{fig:toy_ae} illustrates that optimal changes in model semantics are achieved through adjustments to neuron directions—a process that OFT directly enables. Moreover, the optimization in OFT can be interpreted as tuning neurons constrained to move along a smooth hypersphere manifold, thereby fostering a more favorable optimization geometry. Experimental observations in \cite{liu2021orthogonal} further corroborate these theoretical advantages.

\textbf{Why not minimize hyperspherical energy}. Unlike \cite{liu2021orthogonal}, our approach does not pursue the minimization of hyperspherical energy. In the context of classification tasks, redundancy-free neurons are important; enforcing minimum hyperspherical energy leads to uniformly distributed neurons on the hypersphere, a property misaligned with effective finetuning and likely to erase crucial pretraining knowledge.

\textbf{Balancing flexibility and regularization in finetuning}. Our study highlights a fundamental tension between the adaptability and structural constraint of finetuning approaches. While traditional finetuning offers the greatest flexibility, it frequently sacrifices model robustness and is prone to instability and degradation. OFT, despite its straightforward design, achieves a well-calibrated compromise between adaptability and structured regularity by maintaining pairwise angular relationships among neurons. Moreover, implementing block-diagonal parameterization serves as a stricter regularization mechanism for the orthogonal transformation.

\textbf{No extra computational cost during inference}. In contrast to ControlNet, the OFT approach does not increase the inference time or resource requirements of the adapted model. During inference, it suffices to integrate the optimized orthogonal matrix $\bm{R}$ with the original pretrained weight matrix $\bm{W}^0$, forming a new weight matrix $\thickmuskip=2mu \medmuskip=2mu \bm{W}=\bm{R}\bm{W}^0$. This preserves the inference efficiency found in the original pretrained model.

\section{Experiments and Results}

\textbf{Overall experimental protocol}. Our empirical analysis utilizes Stable Diffusion v1.5~\cite{rombach2022high} as the foundation for text-to-image generation. To maintain fairness, image outputs for each method are sampled randomly. For subject-specific image generation, our procedure is based largely on DreamBooth~\cite{ruiz2023dreambooth}; for experiments involving controlled generation, we employ protocols similar to those from ControlNet~\cite{zhang2023adding} and T2I-Adapter~\cite{mou2023t2i}. Additionally, in comparisons with LoRA, both OFT and COFT are restricted to the identical layers targeted by LoRA. Further experimental configurations and expanded results are documented in Appendix~\ref{app:settings}.

\subsection{Subject-driven Generation}

\textbf{Experimental configuration}. For this series of experiments, we select DreamBooth~\cite{ruiz2023dreambooth} and LoRA~\cite{hulora2022} as comparison benchmarks. Each technique applies the loss function defined by DreamBooth. The baseline methods are executed according to their respective original papers, employing their recommended hyperparameter choices. Additional outcomes are included in Appendix~\ref{app:settings},~\ref{app:coft_oft},~\ref{app:more_qualitative}, and~\ref{app:failure}.

\textbf{Finetuning stability and convergence}. We begin by assessing the stability and convergence speed during finetuning for DreamBooth, LoRA, OFT, and COFT. These comparisons are illustrated in Figure~\ref{fig:teaser} and Figure~\ref{fig:db_convergence}. From these results, it is evident that both COFT and OFT enable the diffusion model to be finetuned with substantial stability. Within 400 iterations, DreamBooth and both OFT variants exhibit effective control, whereas LoRA does not succeed in maintaining subject identity. By 2000 iterations, DreamBooth demonstrates image collapse, and LoRA produces yellow fur instead of a yellow shirt. Conversely, OFT and COFT continue to provide reliable and robust control over the output images. These observations corroborate the rapid and steady convergence offered by our OFT framework for subject-driven synthesis. It is important to highlight the robustness of OFT and COFT with respect to the number of finetuning iterations—this insensitivity is highly advantageous, as it minimizes the need to precisely tune iteration counts for different subjects. Consequently, one can assign a relatively high number of iterations for both OFT and COFT without the necessity for extensive adjustment. Specifically, with an appropriately chosen $\epsilon$ for COFT, both learning rate and iteration count can be configured with minimal effort.

\textbf{Quantitative comparison}. In accordance with \cite{ruiz2023dreambooth}, we further perform quantitative assessments to gauge subject fidelity (using DINO~\cite{caron2021emerging} and CLIP-I~\cite{radford2021learning}), text prompt fidelity (with CLIP-T~\cite{radford2021learning}), and diversity across samples (evaluated by LPIPS~\cite{zhang2018unreasonable}). CLIP-I reflects the mean cosine similarity of CLIP embeddings computed between real and synthesized images. Similarly, DINO evaluates this similarity with ViT S/16 DINO embeddings replacing those from CLIP-I. CLIP-T measures the average cosine similarity between text prompt and generated image embeddings, as provided by CLIP. Additionally, we report the average LPIPS cosine similarity across samples generated for the same subject using identical prompts. Table~\ref{table:dreambooth_final} demonstrates that both COFT and OFT significantly surpass DreamBooth and LoRA in terms of DINO and CLIP-I metrics, and offer marginally superior or equivalent performance concerning prompt fidelity and image diversity. For a robust evaluation, each method is finetuned on the same pre-trained model across 30 unique random seeds to control for variance. Overall, our findings indicate that OFT not only facilitates superior convergence and stability during training but also achieves consistently higher final performance.

\textbf{Qualitative comparison}. To provide a clearer understanding of the advantages brought by OFT, we present several randomly selected cases of subject-driven image synthesis in Figure~\ref{fig:dreambooth_final}. For consistency and fairness, all the samples shown are generated from identically finetuned models with the various approaches applied, without any separate optimization of the text prompt tailored to the final outputs. In every approach, we report results from the checkpoint that achieves the highest validation CLIP score. Inspection of Figure~\ref{fig:dreambooth_final} reveals that OFT and COFT both maintain strong semantic consistency with respect to subject identity, whereas LoRA frequently does not retain this identity (for instance, in the bowl scenario, LoRA fails to depict the subject correctly). Additionally, OFT and COFT provide much finer alignment to the conditioning text, whereas DreamBooth, though effective in subject identity preservation, often generates results that do not match the prompt (as illustrated in the first row of the bowl example). Overall, this qualitative evaluation highlights that OFT excels in achieving both fine-grained controllability and subject identity fidelity. Furthermore, OFT shows robustness with respect to the number of iterations and continues to perform well even as this increases—a property not shared by DreamBooth and LoRA. A broader set of qualitative results is presented in Appendix~\ref{app:more_qualitative}. Human subject evaluations, reported in Appendix~\ref{app:human}, offer additional evidence supporting OFT's superior performance.

\subsection{Controllable Generation}

\textbf{Settings}. As baseline methods, we employ ControlNet~\cite{zhang2023adding}, T2I-Adapter~\cite{mou2023t2i}, and LoRA~\cite{hulora2022}. The main manuscript focuses on three difficult controllable image generation scenarios: Canny edge to image~(C2I) on the COCO dataset~\cite{lin2014microsoft}, segmentation map to image~(S2I) on the ADE20K dataset~\cite{zhou2017scene}, and landmark to face~(L2F) on the CelebA-HQ dataset~\cite{karras2018progressive,xia2021tedigan}. All models are fine-tuned using Stable Diffusion~(SD) v1.5 across each dataset for 20 training epochs. Additional results and analyses can be found in Appendix~\ref{app:more_qualitative},~\ref{app:more_control_task},~and~\ref{app:failure}.

\textbf{Convergence}. We assess how quickly ControlNet, T2I-Adapter, LoRA, and COFT converge on the L2F task, providing both numerical and visual analyses. As our quantitative metric, we calculate the average $\ell_2$ distance between the facial landmarks used for control and the landmarks predicted by the model. Figure~\ref{fig:face_convergence} presents the landmark error across models fine-tuned for varying numbers of epochs. Notably, both COFT and OFT demonstrate much faster convergence compared to alternatives. For example, LoRA requires 20 epochs to match the performance that OFT reaches after only 8 epochs. Importantly, both OFT and COFT utilize a number of trainable parameters comparable to LoRA (and substantially fewer than ControlNet), yet they converge much more efficiently than these other approaches. The superior convergence behavior of OFT is further supported by results depicted in Figure~\ref{fig:teaser}, where the rightmost illustration indicates OFT’s high data efficiency—achieving satisfactory convergence with just 5\% of ADE20K data, outperforming both ControlNet and LoRA in this regard. For our qualitative analysis, we focus on a comparative study among OFT, COFT, and ControlNet, as ControlNet’s landmark prediction error is closest to ours. Visualizations in Figure~\ref{fig:face_convergence_qual} reveal that OFT and COFT not only converge reliably, but also produce face poses that progressively better align with the input landmarks. By contrast, ControlNet exhibits a delayed, abrupt onset of controllability after eight epochs, corresponding to the sharp convergence transition documented in \cite{zhang2023adding}. The reliable and gradual convergence of our OFT framework renders it significantly more practical, as its training process is far smoother and more predictable, which in turn simplifies the search for optimal hyperparameters.

\textbf{Quantitative comparison}. We propose a metric for control consistency to quantitatively assess controllable generation methods. This approach involves deriving the control signal from each generated image and measuring its similarity to the original control input. Specifically, for the C2I task, we report both IoU and F1 score. In the context of S2I, we evaluate mean IoU as well as mean and overall accuracy. For L2F, the evaluation uses the mean $\ell_2$ distance calculated between the target and predicted landmarks. Additional information regarding the computation of these metrics can be found in Appendix~\ref{app:settings}. When comparing all approaches, we utilize their optimal hyperparameters. According to Table~\ref{table:control_gen}, both OFT and COFT achieve significantly stronger and more precise control than alternative baselines. Adapter-based methods, such as T2I-Adapter and ControlNet, demonstrate slower convergence and ultimately inferior results. LoRA outperforms ControlNet in the S2I setting, but underperforms on both C2I and L2F. Notably, we find that LoRA's upper-bound performance remains limited, despite extensive hyperparameter optimization. In contrast, the OFT framework's performance continues to improve with increased trainable parameters, indicating no clear performance plateau. It is important to note that our systematic quantitative assessment for controllable generation is an original contribution, enabling precise measurement of control fidelity in finetuned models—subject to the reliability of the employed segmentation or detection models.

\textbf{Qualitative comparison}. We further conduct a qualitative assessment of OFT and COFT in relation to contemporary leading approaches such as ControlNet, T2I-Adapter, and LoRA. As illustrated in the randomly sampled outputs in Figure~\ref{fig:control_final}, OFT and COFT consistently deliver images that are both photorealistic and exhibit precise control fidelity. Within the S2I scenario, LoRA is unable to generate images that conform to the input segmentation, whereas ControlNet, OFT, and COFT effectively direct the synthesis process according to the given structure. Notably, OFT and COFT are distinguished from ControlNet by their ability to produce highly detailed, realistic images and well-structured geometric forms, all while utilizing a much smaller model footprint. For the C2I setting, OFT and COFT can convincingly generate semantically consistent images even from coarse Canny edges, outperforming both T2I-Adapter and LoRA, whose results are notably inferior. Addressing the L2F challenge, our method stands out by yielding the most precise facial pose guidance even in scenarios involving difficult facial orientations. Across all three evaluated tasks, OFT and COFT consistently surpass state-of-the-art baselines in qualitative image outcomes, underscoring the practical strengths of the OFT framework for controllable generation. To supplement these comparisons, additional qualitative samples covering all three control tasks can be found in Appendix~\ref{app:control_exp}. Moreover, Appendix~\ref{app:more_control_task} demonstrates that OFT also generalizes effectively to a broader array of control tasks, such as dense pose to human body, sketch to image, and depth to image translation.

\section{Concluding Remarks and Open Problems}

Building upon the insight that angular relationships between neurons fundamentally influence visual semantics, we introduce orthogonal finetuning—a straightforward and efficient finetuning technique—for controlling text-to-image diffusion models. Our approach addresses two principal applications: subject-driven generation and controllable generation. In comparison to prevailing techniques, OFT achieves enhanced control and finetuning robustness while utilizing a reduced set of finetuning parameters. Additionally, OFT preserves inference efficiency, imposing no extra computational burden during deployment.

Despite these advantages, OFT gives rise to several compelling open questions. One challenge stems from OFT's enforcement of orthogonality through Cayley parametrization, which requires computing a matrix inverse—a process that can constrain the scalability of the method. Although block diagonal parametrization partially mitigates this limitation, devising a differentiable and faster strategy for matrix inversion remains unresolved. Another direction involves the compositional properties of OFT: the orthogonal matrices generated by distinct finetuning instances can be combined multiplicatively, always resulting in another orthogonal matrix. It remains to be determined whether such compositions retain the specialized knowledge from all corresponding downstream tasks. Lastly, while the parameter efficiency of OFT benefits from the block diagonal design, this also induces certain biases and restricts model flexibility. Identifying alternative strategies to enhance parameter efficiency without such trade-offs constitutes an important avenue for future research.

\end{document}